% !TEX program = xelatex
% Written assignments include the question stem and necessary options so that the review is self-contained.
\documentclass[11pt,a4paper,fontset=fandol]{ctexart}

\newcommand{\coursecode}{SC2005}
\newcommand{\coursetitle}{Operating Systems}
\newcommand{\academicyear}{AY2026--27 Semester 1}
\usepackage{needspace}
\usepackage[margin=25mm]{geometry}
\usepackage{amsmath,amssymb,amsthm,bm,mathtools}
\usepackage{booktabs,tabularx}
\usepackage{enumitem}
\usepackage{graphicx}
\usepackage{xcolor}
\usepackage{tikz}
\usepackage{fancyhdr}
\usepackage{hyperref}
\usepackage[nameinlink,noabbrev]{cleveref}
\usetikzlibrary{arrows.meta,calc,fit,positioning}

\definecolor{ntuRed}{HTML}{D71440}
\definecolor{noteBlue}{HTML}{1F4E79}
\definecolor{verifyOrange}{HTML}{A64B00}

\hypersetup{
  colorlinks=true,
  linkcolor=noteBlue,
  citecolor=noteBlue,
  urlcolor=noteBlue,
  pdftitle={\coursecode: \coursetitle},
  pdfsubject={\academicyear\ 课程学习笔记}
}

\setlength{\parindent}{0pt}
\setlength{\parskip}{0.6em}
\setlist{itemsep=0.25em,topsep=0.35em}
\graphicspath{{figures/}}

\pagestyle{fancy}
\fancyhf{}
\lhead{\coursecode}
\rhead{\coursetitle}
\cfoot{\thepage}
\setlength{\headheight}{14pt}

\newtheorem{definition}{Definition（定义）}[section]
\newtheorem{theorem}[definition]{Theorem（定理）}
\newtheorem{proposition}[definition]{Proposition（命题）}
\newtheorem{example}[definition]{Example（例）}
\newtheorem{remark}[definition]{Remark（说明）}

\newcommand{\recordid}[1]{%
  \textcolor{noteBlue}{\textbf{编号：} \texttt{#1}}%
}

\newcommand{\lecturemeta}[6]{%
  \begin{center}
    \begin{tabularx}{\textwidth}{@{}lX@{}}
      \toprule
      课堂编号 & \texttt{#1} \\
      上课日期 & #2 \\
      主要资料 & #3 \\
      覆盖范围 & #4 \\
      视频补充 & #5 \\
      人工确认 & #6 \\
      \bottomrule
    \end{tabularx}
  \end{center}
}

\newcommand{\exercisemeta}[5]{%
  \begin{center}
    \begin{tabularx}{\textwidth}{@{}lX@{}}
      \toprule
      题组编号 & \texttt{#1} \\
      对应日期 & #2 \\
      主要来源 & #3 \\
      核对状态 & #4 \\
      人工确认 & #5 \\
      \bottomrule
    \end{tabularx}
  \end{center}
}

\newcommand{\lecturetopic}[2]{%
  \subsection{#2}\label{topic:#1}%
  \recordid{#1}\par
}

\newcommand{\exercisequestion}[2]{%
  \subsection{#2}\label{exercise:#1}%
  \recordid{#1}\par
}

\newcommand{\relatedtopic}[2]{%
  \hyperref[topic:#1]{#2}~(\texttt{#1})%
}

\newcommand{\relatedexercise}[2]{%
  \hyperref[exercise:#1]{#2}~(\texttt{#1})%
}

\newcommand{\codeassignmentref}[2]{%
  \texttt{#1}~(\nolinkurl{#2})%
}

\newcommand{\verify}[1]{%
  \textcolor{verifyOrange}{\textbf{[Verify:} #1\textbf{]}}%
}

\newcommand{\keyidea}[1]{%
  \par\textcolor{noteBlue}{\textbf{核心。}} #1\par
}

\newenvironment{workedsolution}
  {\par\textbf{参考解答。}\par}
  {\hfill$\square$\par}

\newenvironment{checkedsolution}
  {\par\textbf{核对后的解答。}\par}
  {\hfill$\square$\par}

\title{\textcolor{ntuRed}{\coursecode}\\[0.4em]\coursetitle}
\author{课程学习笔记}
\date{\academicyear}

\title{\textcolor{ntuRed}{\coursecode}\\[0.4em]\coursetitle\\[0.8em]\large 书面作业}
\author{作业整理}

\begin{document}
\maketitle

\section{Lecture Quiz 1}
\label{homework:SC2005-HW01}

\exercisemeta
  {SC2005-HW01}
  {2026-08-23}
  {NTULearn Lecture Quiz 1，10 道 multiple-choice questions（选择题）；题干与选项见下文}
  {已完成，并根据平台反馈逐题核对}
  {是（2026-08-23）}

\Needspace{12\baselineskip}
\exercisequestion{SC2005-HW01-Q01}{Question 1：Thread Shared Resources}

\textbf{题干：} What parts of the process memory space are shared between threads in the same process?

\begin{enumerate}[label=\Alph*.,leftmargin=2.4em]
  \item Everything is shared between threads in a single process.
  \item Only Text and OS resources (files, etc.) are shared; all other parts are private to the threads.
  \item Text, Data, Heap and OS resources (files, etc.) are shared between the threads; Stack and Registers are not.
  \item Text, Data and Heap are shared between the threads; OS resources (files, etc.), Stack and Registers are not.
\end{enumerate}

\textbf{正确答案：C}\qquad
\textbf{对应讲义：}\texttt{SC2005-L04-T04}

\Needspace{5\baselineskip}
\begin{checkedsolution}
同一 process（进程）的 threads（线程）共享 text、data、heap 与 open files 等 OS resources；stack、registers 与 program counter 由各 thread 独享。
\end{checkedsolution}

\Needspace{12\baselineskip}
\exercisequestion{SC2005-HW01-Q02}{Question 2：Definition of an Operating System}

\textbf{题干：} What is an Operating System?

\begin{enumerate}[label=\Alph*.,leftmargin=2.4em]
  \item Software program that controls other user programs.
  \item Allocates and manages hardware resources.
  \item Always ready to accept new commands from users and hardware.
  \item All the other options are correct.
\end{enumerate}

\textbf{正确答案：D}\qquad
\textbf{对应讲义：}\texttt{SC2005-L02-T01}

\Needspace{5\baselineskip}
\begin{checkedsolution}
OS 同时承担 control program（控制程序）、resource allocator（资源分配器）和 event-driven manager（事件驱动管理者）的职责，因此前三项共同描述了课程中的 OS 角色。
\end{checkedsolution}

\Needspace{12\baselineskip}
\exercisequestion{SC2005-HW01-Q03}{Question 3：Timer Interrupt}

\textbf{题干：} In the process state transition diagram, what is a timer interrupt?

\begin{enumerate}[label=\Alph*.,leftmargin=2.4em]
  \item It is a hardware interrupt generated by a hardware/software clock that tracks progress of physical time.
  \item It is an interrupt generated by the OS to denote process completion.
  \item It is a trap interrupt generated by the CPU.
  \item None of the options are correct.
\end{enumerate}

\textbf{正确答案：A}\qquad
\textbf{对应讲义：}\texttt{SC2005-L02-T03}，\texttt{SC2005-L03-T02}

\Needspace{5\baselineskip}
\begin{checkedsolution}
Timer interrupt（时钟中断）由 timer hardware（定时器硬件）异步产生，使 OS 重新获得控制权；它不是当前程序主动触发的 trap（陷阱）。
\end{checkedsolution}

\Needspace{12\baselineskip}
\exercisequestion{SC2005-HW01-Q04}{Question 4：Multiprogramming and Parallelism}

\textbf{题干：} Multiprogrammed systems allow parallel execution of jobs, hence they are also called time-shared systems.

\begin{enumerate}[label=\Alph*.,leftmargin=2.4em]
  \item False: Parallel execution of jobs depends on the hardware (multiprocessing); multiprogramming allows efficient switching between jobs for CPU access (time-division access), and hence it is also called a time-shared system.
  \item True: Multiprogramming allows simultaneous access to the CPU for multiple user jobs, and hence it is also called a time-shared system.
\end{enumerate}

\textbf{正确答案：A}\qquad
\textbf{对应讲义：}\texttt{SC2005-L02-T02}

\Needspace{5\baselineskip}
\begin{checkedsolution}
Multiprogramming（多道程序设计）通过切换形成 concurrency（并发）；true parallelism（真正并行）取决于 multicore/multiprocessing hardware（多核/多处理硬件）。更严格地说，time-sharing（分时）是强调 preemption（抢占）与 interactive response（交互响应）的 multiprogramming 扩展。
\end{checkedsolution}

\Needspace{12\baselineskip}
\exercisequestion{SC2005-HW01-Q05}{Question 5：Real-Time Requirement}

\textbf{题干：} In real-time systems, average response time of user programs should be minimized.

\begin{enumerate}[label=\Alph*.,leftmargin=2.4em]
  \item True: This will ensure that the required deadlines are met.
  \item True: This will ensure that the system is responsive.
  \item False: Real-time systems require guaranteed satisfaction of deadlines, irrespective of average response time.
  \item False: Real-time systems require average response times to be maximized.
\end{enumerate}

\textbf{正确答案：C}\qquad
\textbf{对应讲义：}\texttt{SC2005-L02-T02}

\Needspace{5\baselineskip}
\begin{checkedsolution}
Real-time system（实时系统）要求对 deadline（截止期限）给出保证；较小的 average response time（平均响应时间）本身不能提供该保证。
\end{checkedsolution}

\Needspace{12\baselineskip}
\exercisequestion{SC2005-HW01-Q06}{Question 6：Multithreading、Multiprogramming and Multiprocessing}

\textbf{题干：} What is the difference between multi-threading, multi-programming and multi-processing?

\begin{enumerate}[label=\Alph*.,leftmargin=2.4em]
  \item They are all different terms for the same concept of time-multiplexed execution.
  \item Multi-threading and multi-programming are identical; multi-processing means multiple CPU cores for execution.
  \item Multi-threading means multiple threads of execution within a process; multi-programming and multi-processing are identical.
  \item Multiple threads of execution within a process; multiple processes ready for execution in main memory; multiple CPU cores for execution.
\end{enumerate}

\textbf{正确答案：D}\qquad
\textbf{对应讲义：}\texttt{SC2005-L02-T02}，\texttt{SC2005-L04-T04}

\Needspace{5\baselineskip}
\begin{checkedsolution}
Multithreading（多线程）指一个 process 内有多条执行流；multiprogramming 指 memory 中有多个可调度 processes；multiprocessing 指存在多个 CPU cores。
\end{checkedsolution}

\Needspace{12\baselineskip}
\exercisequestion{SC2005-HW01-Q07}{Question 7：One-to-One Thread Mapping}

\textbf{题干：} What is an advantage and a disadvantage of one-to-one mapping of logical (user) and physical (kernel) threads?

\begin{enumerate}[label=\Alph*.,leftmargin=2.4em]
  \item More concurrency than many-to-one mappings; mapping is very complex.
  \item More concurrency than many-to-many mappings; requires creation of several kernel threads (one per user thread).
  \item More concurrency than many-to-one mappings; requires creation of several kernel threads (one per user thread).
  \item Mapping is very simple; less concurrency than many-to-one mappings.
\end{enumerate}

\textbf{正确答案：C}\qquad
\textbf{对应讲义：}\texttt{SC2005-L04-T05}

\Needspace{5\baselineskip}
\begin{checkedsolution}
One-to-One mapping（一对一映射）比 Many-to-One 提供更高 concurrency，并允许多核并行；代价是每个 user thread 都需要一个 kernel thread。
\end{checkedsolution}

\Needspace{12\baselineskip}
\exercisequestion{SC2005-HW01-Q08}{Question 8：Privileged I/O Instruction}

\textbf{题干：} When a user program tries to execute an I/O instruction directly,

\begin{enumerate}[label=\Alph*.,leftmargin=2.4em]
  \item the CPU executes the instruction.
  \item the CPU checks if the request is within range of the base and limit registers.
  \item the CPU checks for its validity and legality before executing it.
  \item the CPU generates a trap and transfers control to the OS.
\end{enumerate}

\textbf{正确答案：D}\qquad
\textbf{对应讲义：}\texttt{SC2005-L02-T04}

\Needspace{5\baselineskip}
\begin{checkedsolution}
I/O instruction（输入输出指令）是 privileged instruction（特权指令）；user mode 直接执行会产生 trap，并把控制权交给 OS。
\end{checkedsolution}

\Needspace{12\baselineskip}
\exercisequestion{SC2005-HW01-Q09}{Question 9：Direct and Indirect Message Passing}

\textbf{题干：} How do you distinguish between direct and indirect message passing for inter-process communication?

\begin{enumerate}[label=\Alph*.,leftmargin=2.4em]
  \item In direct message passing, communication occurs between specific processes; in indirect message passing, any process with access to a mailbox can send/receive messages to/from that mailbox.
  \item In direct message passing, communication occurs between specific processes through shared user-space memory; in indirect message passing, communication occurs through mailboxes.
  \item Direct message passing is the same as shared memory; indirect message passing uses the kernel memory space.
\end{enumerate}

\textbf{正确答案：A}\qquad
\textbf{对应讲义：}\texttt{SC2005-L04-T03}

\Needspace{5\baselineskip}
\begin{checkedsolution}
Direct message passing（直接消息传递）显式指定通信 process；indirect message passing（间接消息传递）通过可由多个 processes 访问的 mailbox/port。
\end{checkedsolution}

\Needspace{12\baselineskip}
\exercisequestion{SC2005-HW01-Q10}{Question 10：Process Context and PCB}

\textbf{题干：} What comprises a process context and where is it stored?

\begin{enumerate}[label=\Alph*.,leftmargin=2.4em]
  \item All registers required to execute/resume the process (program counter, stack pointer, memory protection registers, etc.) are stored in the PCB.
  \item The process memory, as defined by the base and limit register values, is stored in the user memory space.
  \item All registers required to resume the process (program counter, stack pointer, memory protection registers, etc.) are stored in the user memory space.
  \item Only base and limit register values for a process are stored in the PCB.
\end{enumerate}

\textbf{正确答案：A}\qquad
\textbf{对应讲义：}\texttt{SC2005-L03-T03}

\Needspace{5\baselineskip}
\begin{checkedsolution}
Process context（进程上下文）包含恢复执行所需的 program counter、stack pointer、registers 与 protection state（保护状态），由 OS 保存在 PCB。
\end{checkedsolution}

\subsection{易错点}

\begin{itemize}
  \item Concurrency（并发）不等于 parallelism（并行）。Time-sharing（分时）更严格地说是强调 preemption（抢占）和 interactive response（交互响应）的 multiprogramming 扩展。
  \item Timer interrupt 属于 asynchronous hardware interrupt（异步硬件中断）；非法 privileged instruction 或 system call 引发的 trap 属于 synchronous event（同步事件）。
  \item Thread 的 shared resources（共享资源）与 private execution context（私有执行上下文）必须分开记忆：heap 共享，stack 独立。
\end{itemize}

\section{Lecture Quiz 2}
\label{homework:SC2005-HW02}

\exercisemeta
  {SC2005-HW02}
  {2026-08-29}
  {NTULearn Lecture Quiz 2，5 道 multiple-choice questions（选择题）；题干与选项见下文}
  {已完成，并根据平台满分反馈逐题核对}
  {是（2026-08-29）}

\Needspace{13\baselineskip}
\exercisequestion{SC2005-HW02-Q01}{Question 1：Preemptive Scheduling Events}

\textbf{题干：} In a preemptive CPU scheduler, when does scheduling happen?

\begin{enumerate}[label=\Alph*.,leftmargin=2.4em]
  \item Upon any of the five transitions (in the process state transition diagram).
  \item Upon any of the five transitions (in the process state transition diagram) and even at other time instants.
  \item Upon transitions 1 and 5 and ocassionally upon transition 4 (in the process state transition diagram).
  \item Upon transitions 2, 3 and 4 (in the process state transition diagram).
\end{enumerate}

\textbf{正确答案：B}\qquad
\textbf{对应讲义：}\texttt{SC2005-L05-T02}

\Needspace{6\baselineskip}
\begin{checkedsolution}
Preemptive scheduling（抢占式调度）不仅包含五种 state transitions（状态转换）上的调度，还允许 scheduler 在其他 time instants（时刻）收回 CPU，例如 time quantum（时间片）到期。D 只列出 transitions 2、3、4，遗漏了 transitions 1、5 以及其他时刻。Scheduler 被唤起并重新选择不等于一定发生 process context switch（进程上下文切换）。
\end{checkedsolution}

\Needspace{12\baselineskip}
\exercisequestion{SC2005-HW02-Q02}{Question 2：Aging}

\textbf{题干：} What is Aging?

\begin{enumerate}[label=\Alph*.,leftmargin=2.4em]
  \item A technique in which the priority of processes that are unable to execute is slowly increased over time to avoid starvation.
  \item A technique in which the priority of processes that are unable to execute is slowly decreased over time to avoid starvation.
  \item A technique in which the priority of all processes is slowly increased over time to avoid starvation.
\end{enumerate}

\textbf{正确答案：A}\qquad
\textbf{对应讲义：}\texttt{SC2005-L06-T02}

\Needspace{6\baselineskip}
\begin{checkedsolution}
Aging（老化）逐步提高长期无法执行的 processes 的 effective priority（有效优先级），使其最终获得 CPU，从而避免 starvation（饥饿）。这里的“提高 priority”指提高调度地位；若系统采用 smaller integer means higher priority（数字越小优先级越高），实现时反而会逐渐减小 priority number。
\end{checkedsolution}

\Needspace{13\baselineskip}
\exercisequestion{SC2005-HW02-Q03}{Question 3：Global Multiprocessor Scheduling}

\textbf{题干：} Under global multiprocessor scheduling, a process may execute on two cores at the same time.

\begin{enumerate}[label=\Alph*.,leftmargin=2.4em]
  \item False, a process can only execute on one core at any time. Under global scheduling, it may migrate between cores over time.
  \item False, a process can only execute on one core at any time. Under global scheduling, a process is assigned to a core when it arrives in the system and is always executed only on that core.
  \item True, under global scheduling, a process may migrate between cores over time and execute on them in parallel.
  \item True, under global scheduling, a process may execute on two cores in parallel, but never on three or more cores in parallel.
\end{enumerate}

\textbf{正确答案：A}\qquad
\textbf{对应讲义：}\texttt{SC2005-L06-T05}

\Needspace{6\baselineskip}
\begin{checkedsolution}
按本讲模型，一个 process 在任一时刻只能占用一个 CPU core。Global scheduling（全局调度）允许该 process 在不同调度区间运行于不同 cores，即发生 migration（迁移），但不允许同一个 execution stream 同时在多个 cores 上执行。B 描述的是固定 core mapping（核心映射），属于 partitioned scheduling（分区调度）。
\end{checkedsolution}

\Needspace{13\baselineskip}
\exercisequestion{SC2005-HW02-Q04}{Question 4：RR Quantum 与 Turnaround Time}

\textbf{题干：} Turnaround time of processes decreases as the quantum size in Round-Robin (RR) scheduling increases.

\begin{enumerate}[label=\Alph*.,leftmargin=2.4em]
  \item False, there is no direct correlation between turnaround time and quantum size.
  \item False, it increases as the quantum size increases.
  \item False, it is always fixed independent of the quantum size.
  \item True, and further it remains fixed once the quantum size is larger than the maximum CPU burst length.
\end{enumerate}

\textbf{正确答案：A}\qquad
\textbf{对应讲义：}\texttt{SC2005-L06-T03}

\Needspace{6\baselineskip}
\begin{checkedsolution}
Average turnaround time（平均周转时间）与 time quantum（时间片）之间通常不存在单调关系；改变 $q$ 会改变 completion order（完成顺序），因此可能增大、减小或保持不变。在讲义忽略 context-switch overhead（上下文切换开销）的固定 workload 模型中，当 $q$ 大于最大 CPU burst length 后，RR 退化为 FCFS，turnaround time 才不再随 $q$ 改变。D 的后半句成立，但题干声称“随 $q$ 增大而减小”，所以整体仍为 False。
\end{checkedsolution}

\Needspace{13\baselineskip}
\exercisequestion{SC2005-HW02-Q05}{Question 5：Nonpreemptive Scheduling Events}

\textbf{题干：} In a nonpreemptive CPU scheduler, when does scheduling happen?

\begin{enumerate}[label=\Alph*.,leftmargin=2.4em]
  \item Upon transitions 1 and 5 and ocassionally upon transition 4 (in the process state transition diagram).
  \item Upon transitions 2, 3 and 4 (in the process state transition diagram).
  \item Upon any of the five transitions (in the process state transition diagram).
  \item Upon any of the five transitions (in the process state transition diagram) and even at other time instants.
\end{enumerate}

\textbf{正确答案：A}\qquad
\textbf{对应讲义：}\texttt{SC2005-L05-T02}

\Needspace{6\baselineskip}
\begin{checkedsolution}
Nonpreemptive scheduling（非抢占式调度）只在当前 process 主动释放 CPU 时重新选择：transition 1 为 Running $\rightarrow$ Waiting，transition 5 为 Running $\rightarrow$ Terminated。Transition 4（New $\rightarrow$ Ready）通常不会打断正在运行的 process；只有 CPU core 原本 idle 时才需要立即调度，因此是 occasionally（偶尔）。
\end{checkedsolution}

\subsection{易错点}

\begin{itemize}
  \item Scheduling decision（调度决策）、process preemption（进程抢占）和 process context switch 是三个不同事件；scheduler 运行后仍可能选择原 process。
  \item Aging 提高的是 effective priority，而不保证 priority number 数值增大；必须先确认题目的优先级编号约定。
  \item Global scheduling 允许同一 process 跨 cores 迁移，不表示同一个 process 能在多个 cores 上同时执行。
  \item RR 的 $q$ 与 turnaround time 没有一般性的单调关系；$q$ 足够大时退化为 FCFS，只是一个边界情况。
\end{itemize}

\section{Lecture Quiz 3}
\label{homework:SC2005-HW03}

\exercisemeta
  {SC2005-HW03}
  {2026-09-12}
  {NTULearn Lecture Quiz 3，10 道 multiple-choice questions（选择题）；题干与选项见下文}
  {已完成，并根据平台满分反馈逐题核对}
  {是（2026-09-12）}

\Needspace{13\baselineskip}
\exercisequestion{SC2005-HW03-Q01}{Question 1：Mutual Exclusion}

\textbf{题干：} What is meant by the mutual exclusion property?

\begin{enumerate}[label=\Alph*.,leftmargin=2.4em]
  \item No two processes can be in the critical section at the same time.
  \item At most two processes can be in the critical section at the same time.
  \item No process can context switch in the critical section.
  \item The OS cannot context switch from one process in its critical section to another process in its critical section.
\end{enumerate}

\textbf{正确答案：A}\qquad
\textbf{对应讲义：}\texttt{SC2005-L07-T04}

\Needspace{6\baselineskip}
\begin{checkedsolution}
Mutual exclusion（互斥）要求任意时刻至多一个 process（进程）处于相关 critical section（临界区）。它限制的是并发访问结果，并不禁止 critical section 内发生 context switch（上下文切换）。
\end{checkedsolution}

\Needspace{13\baselineskip}
\exercisequestion{SC2005-HW03-Q02}{Question 2：Blocking Wait 的操作顺序}

\textbf{题干：} In the blocking implementation of \texttt{Wait()}, why can we not decrement the semaphore value after the \texttt{if()} condition check?

\begin{enumerate}[label=\Alph*.,leftmargin=2.4em]
  \item This would imply a context-switch between the \texttt{if()} condition check and the decrement, which can lead to mutual exclusion violation.
  \item This would imply a context-switch between the \texttt{if()} condition check and the decrement, which can lead to a race condition for the semaphore value.
  \item It is possible, but the implementation is less efficient than when the decrement is before the \texttt{if()} condition check.
\end{enumerate}

\textbf{正确答案：A}\qquad
\textbf{对应讲义：}\texttt{SC2005-L09-T04}，\texttt{SC2005-L09-T05}

\Needspace{8\baselineskip}
\begin{checkedsolution}
若先检查再 decrement（减一），当 $S=1$ 时，P0 可通过检查后被切走；P1 也通过检查、减一并进入 critical section，随后 P0 恢复、减一并同样进入。Check--update（检查—更新）不是一个 atomic action（原子操作），因而可能直接破坏 mutual exclusion。正确顺序是先执行 $S.value\mathrel{-}=1$，再按更新后的值决定是否 block（阻塞）。
\end{checkedsolution}

\Needspace{14\baselineskip}
\exercisequestion{SC2005-HW03-Q03}{Question 3：Blocking Semaphore 的 Atomicity}

\textbf{题干：} Why must \texttt{Wait()} and \texttt{Signal()} system calls with blocking implementation be atomic? Give all reasons.

\begin{enumerate}[label=\Alph*.,leftmargin=2.4em]
  \item To prevent race condition in the updates to the value of the semaphore; To ensure mutual exclusion by guaranteeing that the update to the semaphore value and its checking in the \texttt{if()} condition happen atomically in each system call.
  \item To ensure mutual exclusion by guaranteeing that the update to the semaphore value and its checking in the \texttt{if()} condition happen atomically in each system call.
  \item To prevent race condition in the updates to the value of the semaphore.
  \item They don't have to be atomic. Only busy waiting implementations require atomic \texttt{Wait()} and \texttt{Signal()} system calls.
\end{enumerate}

\textbf{正确答案：A}\qquad
\textbf{对应讲义：}\texttt{SC2005-L09-T05}

\Needspace{7\baselineskip}
\begin{checkedsolution}
Atomicity（原子性）同时保护两件事：多个 callers（调用者）不能对共享 semaphore value（信号量值）产生 lost update（丢失更新）；value update、condition check 以及必要的 queue operation（队列操作）也不能被其他调用插入，否则 mutual exclusion 或 wake-up bookkeeping（唤醒记账）可能失效。
\end{checkedsolution}

\Needspace{13\baselineskip}
\exercisequestion{SC2005-HW03-Q04}{Question 4：Algorithm 2 的 Progress Violation}

\textbf{题干：} In Algorithm 2 in the lecture, Progress property is violated because

\begin{enumerate}[label=\Alph*.,leftmargin=2.4em]
  \item Process P0 updates flag to true, context switch to Process P1, Process P1 updates flag to true, both the processes are now stuck in the entry while loop.
  \item Process P0 updates flag to false, context switch to Process P1, Process P1 updates flag to false, both the processes are now stuck in the entry while loop.
  \item Process P0 updates flag to true and enters critical section, context switch to Process P1, Process P1 updates flag to true and is now stuck indefinitely.
\end{enumerate}

\textbf{正确答案：A}\qquad
\textbf{对应讲义：}\texttt{SC2005-L08-T02}

\Needspace{7\baselineskip}
\begin{checkedsolution}
Algorithm 2（interest flags）中，两者若先后设置自己的 $flag=true$，随后都会在 \texttt{while (flag[other])} 中等待。此时 critical section 为空，却没有任何 process 能进入，形成 deadlock（死锁）并违反 Progress（进展性）。
\end{checkedsolution}

\Needspace{13\baselineskip}
\exercisequestion{SC2005-HW03-Q05}{Question 5：Progress Property}

\textbf{题干：} What is meant by the Progress property?

\begin{enumerate}[label=\Alph*.,leftmargin=2.4em]
  \item It is liveness; when no process is in the critical section, then a process cannot wait indefinitely to get access to the critical section.
  \item It is fairness; when a process wants access to the critical section, another process cannot repeatedly get access to the critical section indefinitely.
  \item It is liveness; when one process is in the critical section, another process cannot wait indefinitely to get access to the critical section.
\end{enumerate}

\textbf{正确答案：A}\qquad
\textbf{对应讲义：}\texttt{SC2005-L07-T04}

\Needspace{6\baselineskip}
\begin{checkedsolution}
Progress（进展性）是一项 liveness（活性）要求：当 critical section 为空且存在申请者时，选择下一位进入者的决定不能无限延期。它不要求打断当前已在 critical section 中的 process。
\end{checkedsolution}

\Needspace{13\baselineskip}
\exercisequestion{SC2005-HW03-Q06}{Question 6：Reader Preference}

\textbf{题干：} In the Reader-Writer implementation, explain what it means by ``the reader is given preference''.

\begin{enumerate}[label=\Alph*.,leftmargin=2.4em]
  \item While a reader is reading, other readers will be allowed access to the database, even if writers are currently waiting to get access.
  \item If a reader and writer both want access to the database at the same time, then the reader is given preference over the writer irrespective of who is currently accessing the database.
  \item Readers can always access the database, even if writers are currently accessing it.
\end{enumerate}

\textbf{正确答案：A}\qquad
\textbf{对应讲义：}\texttt{SC2005-L10-T05}

\Needspace{7\baselineskip}
\begin{checkedsolution}
Reader preference（读者优先）允许已有 reader group（读者组）存在时，新 readers 继续加入，即使 writer 正在等待；但 reader 与 writer 仍不能同时访问数据库。持续到来的 readers 因此可能造成 writer starvation（写者饥饿）。
\end{checkedsolution}

\Needspace{13\baselineskip}
\exercisequestion{SC2005-HW03-Q07}{Question 7：Bounded Waiting}

\textbf{题干：} What is meant by the property of Bounded Waiting?

\begin{enumerate}[label=\Alph*.,leftmargin=2.4em]
  \item It is fairness; if a process wants to get access to a critical section, another process cannot get access to the critical section repeatedly infinite number of times.
  \item It is liveness; if no process is in the critical section, then a process wanting to get access to the critical section should not wait indefinitely.
  \item It is fairness; processes should make progress collectively through their critical sections.
\end{enumerate}

\textbf{正确答案：A}\qquad
\textbf{对应讲义：}\texttt{SC2005-L07-T04}

\Needspace{6\baselineskip}
\begin{checkedsolution}
Bounded waiting（有界等待）是一项 fairness（公平性）要求：process 发出进入请求后，其他 processes 在它之前进入 critical section 的次数必须有有限上界，从而排除无限 overtaking（超越）和 starvation。
\end{checkedsolution}

\Needspace{15\baselineskip}
\exercisequestion{SC2005-HW03-Q08}{Question 8：Peterson's Algorithm 的 Progress}

\textbf{题干：} Algorithm 3 in the lecture satisfies Progress property because

\begin{enumerate}[label=\Alph*.,leftmargin=2.4em]
  \item Suppose Process P0 wants access to the critical section; then flag for Process P0 is true; if Process P1 is in the entry section then either P0 or P1 will get access to the critical section; if Process P1 is in the critical section then Progress is vacuously satisfied; if Process P1 is in the remainder section then Process P0 can get access to the critical section.
  \item Progress property is satisfied because Bounded Waiting is satisfied.
  \item Process P0 sets its flag to true, turn to 1, and waits at the entry while loop; context switch to Process P1; Process P1 sets its flag to true, turn to 0, and waits at the entry while loop; context switch to Process P0; Process P0 enters the critical section.
\end{enumerate}

\textbf{正确答案：A}\qquad
\textbf{对应讲义：}\texttt{SC2005-L08-T03}

\Needspace{8\baselineskip}
\begin{checkedsolution}
Peterson's Algorithm（Peterson 算法）用 $flag$ 表示意愿、用 $turn$ 在双方同时申请时打破平局：只有一方申请时它可直接进入；双方同时申请时，$turn$ 保证至少一方进入。因此不会在 critical section 为空且有申请者时让两者永久停在 entry section。C 只验证一个 interleaving（交错序列），A 才覆盖 P1 所处的各类状态。
\end{checkedsolution}

\Needspace{14\baselineskip}
\exercisequestion{SC2005-HW03-Q09}{Question 9：Algorithm 2 的 Bounded Waiting}

\textbf{题干：} In Algorithm 2 in the lecture, Bounded Waiting is satisfied because

\begin{enumerate}[label=\Alph*.,leftmargin=2.4em]
  \item A process sets its flag to true at the beginning of entry section; once the flag is set to true another process cannot subsequently enter its critical section.
  \item Progress is not satisfied and this implies that Bounded Waiting will be satisfied.
  \item Process P0 starts and enters critical section, context switch to Process P1, Process P1 sets flag to true and gets stuck in the entry while loop, context switch to Process P0, Process P0 completes critical and remainder sections and tries to enter the critical section again, Process P0 gets stuck.
\end{enumerate}

\textbf{正确答案：A}\qquad
\textbf{对应讲义：}\texttt{SC2005-L08-T02}

\Needspace{8\baselineskip}
\begin{checkedsolution}
一旦 Pi 设置 $flag[i]=true$，Pj 此后不能再次通过自己的 entry loop 进入 critical section，因此 Pj 不可能无限次 overtaking Pi；这满足本讲单独检查的 bounded-waiting 条件。但 Algorithm 2 仍可能让双方一起卡住，所以 bounded waiting 成立并不推出 Progress 成立。C 只是一个特定 trace（执行轨迹），A 给出了通用判断依据。
\end{checkedsolution}

\Needspace{15\baselineskip}
\exercisequestion{SC2005-HW03-Q10}{Question 10：Bounded-Buffer 的 Wait 顺序}

\textbf{题干：} In the Producer-Consumer Bounder-Buffer implementation, why is the \texttt{Wait()} on empty or full semaphore done before the \texttt{Wait()} on mutex?

\begin{enumerate}[label=\Alph*.,leftmargin=2.4em]
  \item Otherwise, it can lead to a deadlock when the buffer is empty/full, because to add/remove items in the buffer the mutex semaphore must be acquired.
  \item Otherwise, it can lead to a deadlock when the buffer is empty/full, because to add/remove items in the buffer the empty/full semaphore must be acquired.
  \item Otherwise, it can lead to a deadlock when the buffer is empty/full, because to add/remove items in the buffer the full/empty semaphore must be acquired.
\end{enumerate}

\textbf{正确答案：A}\qquad
\textbf{对应讲义：}\texttt{SC2005-L10-T03}

\Needspace{9\baselineskip}
\begin{checkedsolution}
若 producer（生产者）先取得 $mutex$，再因 full buffer（满缓冲区）等待 $empty$，consumer（消费者）便无法取得 $mutex$ 来移除 item 并执行 \texttt{Signal(empty)}，双方形成 circular wait（循环等待）；empty buffer 时同理。故应先等待 resource-count semaphore（资源计数信号量）$empty/full$，确认资源可用后再取得 $mutex$ 修改缓冲区。
\end{checkedsolution}

\subsection{易错点}

\begin{itemize}
  \item Mutual exclusion、Progress 与 Bounded Waiting 是三个独立检查项：分别回答“能否同时进入”“空闲时能否继续前进”“能否被无限超越”。
  \item Blocking semaphore 仍要求 atomic \texttt{Wait}/\texttt{Signal}；blocking 只改变等待方式，不会自动消除共享 value 和 waiting queue 上的 race condition。
  \item Algorithm 2 满足本讲的 bounded-waiting 判断，但不满足 Progress；Algorithm 3 通过 $turn$ 解决双方同时声明意愿时的 deadlock。
  \item Producer--consumer 中先等待 $empty/full$、后等待 $mutex$；反序会让 process 持有 $mutex$ 睡眠，使能够改变资源计数的对方也无法进入。
\end{itemize}

\end{document}
